% !TEX root = ../../main.tex

\subsection{Data Representation and Subtree Objects}

\Cref{tab:notation} fixes the global symbols used in the rest of the
method. This subsection introduces the tree objects in plain language before
writing them formally, because the statistical tests operate on subtree
summaries rather than directly on individual samples. All supported inputs are
first mapped into one working sample-feature matrix $X$ with an explicit
feature-space contract. In the current benchmark implementation, Bernoulli and
one-hot cases use Hamming distance and average linkage for the candidate tree;
Gaussian location-style continuous benchmark cases use a precomputed
standardized Euclidean distance matrix, while continuous covariance and
Brownian diagnostic cases may declare separate precomputed tree distances.

For a non-expert reader, the important point is that the tree is a
bookkeeping device for nested groups of samples. Each leaf corresponds to one
observed sample. Each internal node is created by one agglomerative merge and
therefore stands for a group of leaves rather than a new observation. The root
is the unique node that contains every sample. In that sense, internal nodes
represent candidate subpopulations, while leaves represent the original data
points.

The edges of the tree record immediate containment relations between these
groups. If $(u,c) \in E$, then $c$ is a child subtree sitting directly inside
the larger parent subtree $u$. Equivalently, every leaf below $c$ is also
below $u$, and there is no intermediate node between them. This interpretation
is important because the edge test is attached to exactly that parent-child
relation: it asks whether the child subtree distribution departs from the
background defined by its parent. The sibling test then uses the two child
edges under the same parent to ask whether those child subtrees should remain
separate.

In the direct binary or indicator case, $X \in \{0,1\}^{n \times p}$ with rows
$X_i \in \{0,1\}^p$. The leaf set of $\mathcal{T}$ consists of the original
samples, and each agglomerative merge creates one new internal node whose
descendant set is the union of the two merged child subtrees. Starting from $n$
singleton leaves and merging until all samples belong to one connected
hierarchy yields a rooted binary tree with
\begin{equation}
|V| = 2n-1,
\qquad
|E| = 2n-2,
\end{equation}
consisting of $n$ leaves, $n-1$ internal nodes, and one root. The root
represents the full sample set, and each internal node defines one candidate
split location that may or may not survive the later statistical filtering.
Formally, if $u \in V$ is any node, then its associated subtree object is the
set of descendant leaves
\begin{equation}
L(u) = \{ i \in \{1,\dots,n\} : \text{leaf } L_i \text{ lies below } u \}.
\end{equation}
The subtree size is
\begin{equation}
n_u = |L(u)|.
\end{equation}
If $c$ is a child of $u$, then the edge $(u,c)$ represents the immediate set
inclusion
\begin{equation}
L(c) \subset L(u),
\end{equation}
and if $u$ has children $l(u)$ and $r(u)$, then the binary-tree structure gives
the disjoint union identity
\begin{equation}
L(u) = L\bigl(l(u)\bigr) \cup L\bigl(r(u)\bigr),
\qquad
L\bigl(l(u)\bigr) \cap L\bigl(r(u)\bigr) = \varnothing.
\end{equation}

Once that tree is fixed, the inferential problem is about subtrees rather than
individual samples. The method therefore replaces each node by an empirical
feature distribution over its descendant leaves. This is the object that enters
both tests: the edge test compares a child distribution to its parent's
distribution, and the sibling test compares the two child distributions under
the same parent. The raw leaf matrix still matters through the descendant sets
and local spectral matrices, but the primary inferential units are node-level
distributions.

For every supported feature block, the node distribution is the leaf-count
weighted empirical barycenter of the descendant rows. In raw coordinates,
\begin{equation}
\theta_u = \frac{1}{n_u}\sum_{i \in L(u)} X_i.
\end{equation}
Thus $\theta_u$ is not a parameter attached to a single sample. It is the
empirical summary of the whole subtree rooted at $u$: each coordinate
$\theta_{u,j}$ is a descendant mean. For Bernoulli and one-hot coordinates this
mean is a feature or category frequency; for continuous coordinates it is the
empirical subtree mean in the raw continuous coordinate.
If $u$ has children $l$ and $r$, then
\begin{equation}
n_u = n_{l(u)} + n_{r(u)},
\qquad
\theta_u = \frac{n_{l(u)} \theta_{l(u)} + n_{r(u)} \theta_{r(u)}}{n_u}.
\end{equation}
Equivalently, with
\[
\beta_u = \frac{n_{l(u)}}{n_u},
\]
the parent is the linear empirical barycenter
\begin{equation}
\theta_u
=
\beta_u \theta_{l(u)}
+
(1-\beta_u)\theta_{r(u)}.
\label{eq:barycentric-parent}
\end{equation}
This is a centroid representation in the declared feature chart. Euclidean
means are squared-error centroids, and Bregman clustering theory gives the
analogous mean-as-representative interpretation for broad classes of
information-geometric distortions \citep{banerjee2005clustering}. Functional
and symmetrized Bregman centroid results give useful extensions for
distributional summaries \citep{frigyik2008functional,nielsen2009sided}, and
agglomerative Bregman clustering connects this centroid view to hierarchical
merge procedures \citep{telgarsky2012agglomerative}. Ward's minimum-variance
method is the classical Euclidean special case in which merge costs are written
through centroids and cluster sizes \citep{ward1963hierarchical}.

\Cref{eq:barycentric-parent} also gives an exact local algebraic identity for
the edge and sibling tests. For a binary parent with left child $L$ and right
child $R$,
\begin{equation}
\theta_L-\theta_u
=
(1-\beta_u)(\theta_L-\theta_R),
\qquad
\theta_R-\theta_u
=
-\beta_u(\theta_L-\theta_R).
\label{eq:barycentric-raw-contrasts}
\end{equation}
The sibling variance scale is
\[
s_u^{\mathrm{sib}}
=
\left(\frac{1}{n_L}+\frac{1}{n_R}\right)
\left(1+\frac{t_{u,L}+t_{u,R}}{2\bar t}\right),
\]
while the nested left and right child-parent variance scales are
\[
\left(\frac{1}{n_L}-\frac{1}{n_u}\right)
\left(1+\frac{t_{u,L}}{\bar t}\right)
\quad\text{and}\quad
\left(\frac{1}{n_R}-\frac{1}{n_u}\right)
\left(1+\frac{t_{u,R}}{\bar t}\right).
\]
Without branch-time inflation, under a shared local covariance baseline and the
same whitening chart, the barycentric coefficients in
\Cref{eq:barycentric-raw-contrasts} cancel the nested variance scales:
\begin{equation}
z_{L\to u}^{\mathrm{edge}} = z_u^{\mathrm{sib}},
\qquad
z_{R\to u}^{\mathrm{edge}} = -z_u^{\mathrm{sib}}.
\label{eq:barycentric-edge-sibling-identity}
\end{equation}
Thus the same null-whitened barycentric direction can both open the
child-parent edge path and create the focal sibling statistic. This is why the
selected-hierarchy calibration problem is better viewed as a selected
barycentric-contrast problem than as a marginal chi-square calculation.

The barycentric algebra also supplies predeclared leverage variables for
selected-tail diagnostics:
\begin{equation}
b_u=\min(\beta_u,1-\beta_u),
\qquad
\ell_u^{\mathrm{bar}}
=
\log\frac{\max(\beta_u,1-\beta_u)}{\min(\beta_u,1-\beta_u)},
\qquad
s_u^{\mathrm{samp}}=\frac{1}{n_L}+\frac{1}{n_R}.
\end{equation}
In the current diagnostics these variables enter only as descriptive selected
tail coordinates, for example through models of the form
\begin{equation}
\log R_u
=
m\!\left(
A_u^{\mathrm{edge}},
\log n_u,
\log(p/n_u),
k_u^{\mathrm{sib}},
f_u,
b_u,
\ell_u^{\mathrm{bar}},
\lambda_{k,u}/\lambda_{+,u},
M_{k,u},
\rho_u
\right)
+
\varepsilon_u.
\end{equation}
They do not define a production external calibration rule unless held-out tail
validation succeeds in the target selected context.
The decomposition goal is to identify a set of terminal tree nodes whose leaf
sets partition the samples and whose defining splits are supported by the two
statistical tests below. This barycentric subtree update is shared by
Bernoulli, categorical one-hot, and continuous coordinates; the feature-space
contract determines the contrast chart and covariance model used by the tests.

This current representation is still a coordinatewise empirical barycenter, not
a Wasserstein or Fr{\'e}chet barycenter of full probability measures. Future
extensions that represent nodes as full distributions or tree-valued objects
would need their own barycenter geometry: Wasserstein barycenters for
probability measures and Gaussian location-scatter families
\citep{agueh2011barycenters,cuturi2014fast,alvarez2016fixed,carlier2024quantitative},
Fr{\'e}chet or Karcher means for metric and Riemannian spaces
\citep{frechet1948elements,karcher1977riemannian}, and BHV tree-space means
for phylogenetic or topological summaries
\citep{miller2015polyhedral,brown2020mean,lammers2024statistics}.

Categorical features with values $C_{ij} \in \{0,\dots,K-1\}$ are represented
through one-hot indicators
\begin{equation}
\tilde{X}_{i,j,k} = \mathbf{1}\{C_{ij} = k\},
\qquad
k = 0,\dots,K-1,
\end{equation}
so that the raw matrix contains the expanded indicator coordinates. The
statistical contract does not treat those coordinates as independent Bernoulli
variables. Instead, the feature-space object groups each one-hot variable into
a drop-last simplex chart and uses the corresponding multinomial block
covariance in the projected-Wald contrast.
[TODO: validate the effect of one-hot compositional constraints on the
projected edge and sibling tests, especially for high-cardinality categorical
features.]
Continuous inputs can enter through an explicit continuous feature-space
contract. In that contract, the raw continuous columns form an
empirical-Gaussian block, and each tree node stores the empirical descendant
covariance matrix for that block. The active Gaussian location-style continuous
benchmark cases keep the raw continuous values and use a precomputed
standardized Euclidean distance matrix for tree construction; covariance,
low-rank, and Brownian diagnostic cases may instead declare a metric-specific
precomputed distance. Separate discretized Gaussian benchmark variants still
exist for comparison. One such variant is median binarization,
\begin{equation}
X_{ij}^{\mathrm{bin}}
=
\mathbf{1}\!\left\{
W_{ij} > \mathrm{median}(W_{\cdot j})
\right\},
\end{equation}
and another is quantile discretization followed by one-hot expansion,
\begin{equation}
C_{ij}
=
\sum_{q=1}^{K-1} \mathbf{1}\{W_{ij} > t_{jq}\},
\qquad
t_{jq} = Q_{q/K}(W_{\cdot j}),
\end{equation}
with the same indicator map
$\tilde{X}_{i,j,k} = \mathbf{1}\{C_{ij} = k\}$. These discretized variants are
not the continuous-data contract itself. After representation is chosen, the
testing logic reads the node-distribution and feature-space contract; the
tree-distance contract is recorded separately for binary, categorical, and
continuous benchmark cases.
